Figure~\ref{fig:curves} and Table~\ref{tab:main} distinguish fitting from generalization. At decay 1, sustained training interpolation occurs in 0/5 runs at width 16, 5/5 runs at width 32, 5/5 runs at width 64, 5/5 runs at width 128, 5/5 runs at width 256. The corresponding generalization counts are 0/5, 0/5, 0/5, 5/5, 5/5. These observations demonstrate why a single grokking count would conflate insufficient fitting, delayed success, and unresolved generalization.

For the primary paired comparison, widths 64 and 256 have mean restricted generalization waiting times of 29,800 and 11,580 updates, respectively. The width-256 minus width-64 difference is -18,220 [-19,020, -17,180] updates (95\% paired seed bootstrap interval). Restricted training-event waiting times are 300 and 100; their paired difference is -200 [-200, -200] updates. The individual paired outcomes are in the appendix.

Widening therefore changes the slow generalization process much more than the early fitting process at the measured horizon. This contradicts the strong width-independent decay-clock explanation. The narrow model's censored time is a lower bound, so the evidence concerns finite-horizon outcomes and restricted waiting times rather than an estimated eventual event time. It does not establish permanent failure at width 64.

The original five-pair exact two-sided sign test gives $p=0.0625$. Its coarse resolution is a reason not to interpret the bootstrap interval as decisive population-level significance. An independently planned ten-pair replication on seeds 20--29 gives a restricted waiting-time difference of -18,730 [-19,190, -18,260] updates, with 0/10 versus 10/10 events and exact sign-test $p=0.0020$. The replication was designed after seeing the main-grid result and is reported as a separate cohort.

Across the full no-decay arm, 0 of 25 runs reach the generalization criterion. At positive decay, event counts depend on width and decay jointly (Figure~
ef{fig:grid}). In particular, decay 3 gives 0/5 at width 16, 0/5 at width 32, 0/5 at width 64, 0/5 at width 128, 5/5 at width 256. This comparison concerns the same step budget and fixed AdamW learning rate; it does not assert that zero decay precludes grokking under other parameterizations or longer training.

The width-256 decay-clock check reports $\eta\lambda T_{\mathrm{gen}}$, not a rescaled censored time: $\lambda=0.3$: 5/5 events, observed scaled times 17.91--18.99; $\lambda=1$: 5/5 events, observed scaled times 30.90--40.50; $\lambda=3$: 5/5 events, observed scaled times 100.80--157.50.

The ratio of largest to smallest mean scaled event time is 7.62, exceeding the prechosen factor-two descriptive tolerance. This is a diagnostic of approximate collapse, not a population equivalence test.

Among the fully observed reference widths 128 and 256, mean absolute post-fit delays are 17,040 and 11,480 steps, while mean ratios $T_{\mathrm{gen}}/T_{\mathrm{fit}}$ are 86.2 and 115.8. Thus these operational metrics have different width trends. The ratio is especially sensitive to the 100-step evaluation grid and quantization of the short fitting time; it is not evidence of an underlying continuous-time ratio law.

A three-checkpoint event need not remain above threshold. For decay 3 and width 256, 2/5 runs satisfy the additional criterion of staying above 95\% at every remaining observed checkpoint, whereas the original event count is 5/5. At decay 1, the corresponding durable-event count at width 256 is 5/5. These retrospective sensitivity results do not change the prespecified primary endpoint.
